% !TEX program = xelatex
% 宁波市2026学年第一学期高考模拟考试 语文参考答案
\documentclass[11pt,a4paper]{ctexart}
\usepackage{ywexam}

\examfooter{语文参考答案}

\begin{document}

\answertitle{宁波市2026学年第一学期高考模拟考试}{语文参考答案}

\sechead{一、阅读（72分）}

\subsechead{（一）阅读Ⅰ（本题共5小题，19分）}
\ans{1}{D。解析：A项“推动对偶走向成熟”理解错误，汉代“崇尚骈偶”，对偶“更具艺术性”。B项强加因果，“先民二元平衡思维”不是古籍中产生对偶的唯一的、必然的原因。C项，言辞、文辞的概念不是子瑜创设的，他只是引用。}
\ans{2}{C。解析：C项错误。三则材料各承担不同作用，不存在层层递进的说理关系。“言辞与文辞有所区分又一脉相承、不可分割”是作者结合《红楼梦》片段、民间谣谚等案例分析后得出的观点，并非由这三则引文推导得出。}
\ans{3}{B。解析：B项错误，\xq{1}不是反对，只是言对。C项错误，前半句“\xq{3}\xq{5}是言对”正确，但都不是正对。D项错误，\xq{4}确为事对，但上下两句形成正反对照关系，应属反对。}
\ans{4}{\xq{1}甲处：为什么中国如此喜欢用对偶修辞呢？\xq{2}乙处：何以说文字型的文学会有语言的基础呢？}
\scorenote{（每点2分，共4分，写出大致信息即可。）}
\ans{5}{\xq{1}探究这种语言现象的起源和发展历程。\xq{2}对其进行类型划分，评判不同类型的艺术价值。\xq{3}追溯其思想根源，探究其背后的思维习惯、审美心理等。\xq{4}探究这类语言现象的语言基础。}
\scorenote{（每小点2分，共6分，答对3点给满分。）}

\subsechead{（二）阅读Ⅱ（本题共4小题，16分）}
\ans{6}{B。解析：“表现出她对丈夫深切的爱与对‘死书籍’的厌恶”错误，妻子并非厌恶“书籍”。}
\ans{7}{A。解析：“推动情节发展”表述错误。}
\ans{8}{\xq{1}完整呈现书籍写作、排印、售卖、搬运、查扣的流转过程。\xq{2}同时写出社会各阶层的面貌，形成全景式的社会图景。\xq{3}场景之间少有过渡，形成断续跳跃的叙事节奏，留下大量空白，让读者在场景的拼接与对照中深化体会。}
\scorenote{（每点2分，共4分，写出2点即可。）}
\ans{9}{\xq{1}“颜色黯淡”：小说始终笼罩在“黯淡”的色调中——文学者苍白的脸色、昏暗的印刷所与店铺、灰色的街头雪泥等，奠定了作品黯淡的基调，也隐喻着社会环境的冰冷。\xq{2}“铅一般的”：文学者早逝，排字匠怨愤，孩童被书贩、警官等人欺凌，呈现了底层人们“铅一般的”沉重命运。（如果学生从阅读体验的视角，感受到作品给人“铅一般的”沉重的心情也可。）\xq{3}“滑稽”：书籍题为“为了不幸的人们”，却遭排字匠、孩童漠视，还被书贩、警官利用，高尚的理想遭遇“滑稽”的现实，具有反讽的力量。}
\scorenote{（每点3分，共6分，写出2点即可。）}

\subsechead{（三）阅读Ⅲ（本题共5小题，22分）}
\ans{10}{BDG。解析：原文断句为“乃有所窥见于道，而不闻中庸之义，不知所以裁之，遂恣其猖狂而无所极”。}
\ans{11}{A。解析：材料一中的“幸”义为“希望”。}
\ans{12}{B。解析：材料二中的“此知庄子之粗者”是苏轼评价《史记》的话，并非司马迁的观点。}
\ans{13}{（1）认为仆人不爱护公子，是不对的；把这作为侍奉公子的准则，也是不对的。}
\scorenote{（关注“以……为”“以为”“事”“法”的含义。）}
\subans{（2）}{孔子是百代的典范，庄周是什么人，竟敢抨击他，还进而轻慢侮辱他。}
\scorenote{（关注“周”“掊击”“从而”和“狎”的含义。）}
\ans{14}{\xq{1}“文不予”：庄子笔下的孔子见识不高；庄子还借“渔父”之口批评他的做法远离了大道。\xq{2}“实予”：故事展现了孔子心怀天下的仁爱思想、知其不可而为之的进取精神以及谦逊好学的态度，这些都是值得世人学习的美好品质。}
\scorenote{（第\xq{1}点“见识不高”“远离大道”各1分，第\xq{2}点“仁爱思想”“进取精神”和“谦逊好学”各1分，意思对即可。）}

\subsechead{（四）阅读Ⅳ（15分）}
\ans{15}{B。解析：“欣喜”理解错误。颔联表达的是诗人对杜衍突然罢相的惊讶。}
\ans{16}{\xq{1}不恋权位、恬然淡泊。颔联写到他刚刚辞去宰相的职务，便彻底放下为官之念，坦然追求隐居生活。\xq{2}关注国事、心系天下。颈联的“进退忧时”体现了一名儒者的古道热肠和无论在位不在位都心系天下苍生的崇高情怀。\xq{3}低调务实，不慕虚名。尾联说他虽然有着很高的声望，但并不希望自己的功劳事业被谱成乐曲传唱，展现了他低调务实、不慕虚名的品格。}
\scorenote{（每点2分，共6分，意思对即可。）}
\ans{17}{（1）管弦呕哑，多于市人之言语}
\subans{（2）}{长太息以掩涕兮，哀民生之多艰／民生各有所乐兮，余独好修以为常}
\subans{（3）}{示例：舟遥遥以轻飏，风飘飘而吹衣／载舟覆舟，所宜深慎／亲朋无一字，老病有孤舟／沉舟侧畔千帆过，病树前头万木春／只恐双溪舴艋舟，载不动许多愁}
\scorenote{（每空1分，共6分。）}

\sechead{二、语言文字运用（本题共5小题，18分）}
\ans{18}{\xq{1}“秣”改为“抹”；\xq{2}“窜”改为“串”。}
\scorenote{（每处1分，共2分。）}
\ans{19}{D。解析：\xq{3}总起，后面所述内容均为“影片故事”；接下来，先点明时代背景\xq{1}“乱世之中”；再引入人物郑木生\xq{4}；\xq{5}是人物“躲避战乱”的无奈选择，因此紧跟在\xq{4}后；“妻儿留乡”是“远赴南洋”导致的结果，因此放在\xq{5}后。}
\ans{20}{\xq{2}中途易辙，在“大时代”前加上“将”。\xq{3}搭配不当，将“激烈、华丽的冲突和特效”改为“激烈的冲突和华丽的特效”。}
\scorenote{（每处找对序号1分，改错1分，共4分。）}
\ans{21}{\xq{1}“船行入夜，恰江上升明月，圆如玉坠”运用比喻手法，将明月比作玉坠，生动形象，充满画面感。\xq{2}“家中正是收获之季，你切勿操劳，善自珍重”直抒胸臆，朴实话语藏着丈夫远隔千里的牵挂与心疼，真情动人。\xq{3}“暹罗虽远，心有所寄，身若比邻”运用对比，将现实空间的遥远与彼此心灵的紧密连接形成对照，其中专一而笃定的家国之思、亲情之思直击人心。}
\scorenote{（每点2分，共4分。选择其他语句，言之成理亦可给分。）}
\ans{22}{示例：}
\begin{materialtext}
亲爱的兄长：\par
见字如晤！\par
秋意渐浓，窗外满庭桂香，月华如银，清辉洒地，一如当年我们共坐闲谈之夜。而今虽山川阻隔，不能促膝，思念却如长流溪水，涓涓不绝，苏子云：“但愿人长久，千里共婵娟。”望兄适时添衣保暖，善自珍重，并盼早日相聚。\par
\hfill 弟\quad 小安
\end{materialtext}
\scorenote{（引用、抒情和环境描写各1分，语言得体1分，语段完整、有意境1分，共6分。）}

\sechead{三、写作（60分）}
\ans{23}{（60分）作文阅卷参照高考评分要求。}

\par\vspace{0.5em}
\sechead{文言文参考译文}

\translabel{材料一：}
\begin{translation}
孔子在缁帷之林中游玩，坐在杏坛上休息。弟子们在读书，孔子边弹琴边吟唱。曲子还没弹到一半，有位渔父下船走来，左手按着膝盖，右手托着下巴听琴。曲子弹完，他招手叫子贡、子路过去，两人一同上前应答。

那客人指着孔子问：“他是做什么的？”子路回答说：“是鲁国的君子。”客人问他的姓氏宗族。子路说：“孔氏一族。”客人又问：“孔氏研习什么？”子贡回答说：“孔氏本性秉持忠信，亲身践行仁义，打算以此造福天下。这就是孔氏所从事的事业。”客人又问：“他是有封地的君主吗？”子贡说：“不是。”“是侯王的辅佐之臣吗？”子贡说：“不是。”

客人于是笑着往回走，一边笑一边说：“仁倒是仁啊，恐怕不能保全自身；使心神受苦、形体劳累，从而危害了自己的本真。唉，他背离道也太远了！”

子贡回来，把经过告诉孔子。孔子推开琴站起身说：“那大概是位圣人吧！”于是走下杏坛去追他，赶到水边，渔父正拿着船篙要撑开船。客人问：“你要找什么？”孔子说：“刚才先生话还没说完就走了，我孔丘不才，不明白您说的意思，私下在这里恭敬等候，希望能听到您的一言半语来帮助我！”客人说：“嘻！你也太好学了！”孔子拜了两拜才起身，说：“我从少年时就修习学问，直到今天，已经六十九年了，还没能听到至高无上的教诲，怎敢不虚心！”
\end{translation}

\translabel{材料二：}
\begin{translation}
严格按照《史记》的记载来看，庄子与梁惠王、齐宣王是同时代人，他的学问没有什么不涉猎的，然而要旨归属于老子的学说。所以他著书十余万字，大多都是寓言，写下《渔父》《盗跖》《胠箧》等篇，用来非议孔子这类人，来阐明老子的道术。这只是了解庄子的粗浅之处罢了。

我认为庄子大概是帮助孔子的人，关键是不能把他的做法当作准则罢了。楚国公子改换平民服装逃亡，守门人刁难他。他的仆从拿着鞭子骂道：“你这奴才怎么不尽力！”守门人便放他们出去了。事情本来就有看似倒行逆施、实则合乎情理的。认为仆人不爱护公子，是不对的；把这作为侍奉公子的准则，也是不对的。所以庄子的话，都是实际上赞许孔子，而文辞上不赞许；表面上排挤，暗地里帮助，他正面直说的话大概没有多少。
\end{translation}

\translabel{材料三：}
\begin{translation}
唉！孔子是百代的准则典范，庄周是什么人，竟敢抨击他，还进而轻慢侮辱他。不幸他的书广为流传，世上那些喜欢放纵恣肆、害怕拘束检点的人，没有不拿庄子来作借口的，于是导致礼义衰颓，最终使人的家国倾覆败亡，不也是很可悲的吗！
\end{translation}

\translabel{材料四：}
\begin{translation}
韩愈说庄周的学说出自子夏，大概是这样吧。庄周继承了孔氏的末流余绪，于是对道有所窥见领悟，却没有听闻中庸的大义，不懂得如何取舍裁断，便放纵其猖狂恣肆而没有限度。他的文辞如此不谦逊，可苏轼却说他推尊圣人，自居于不能周全、不能通达的、偏于一隅的读书人行列。苏轼对于庄子的理解，相差多么远啊！
\end{translation}

\end{document}
